\section{引言}

近年来，基于大语言模型的自动化分析系统已经能够较稳定地完成代码生成、执行纠错和结果整理等任务，但在更上游的“分析规划”阶段仍普遍存在短板\citep{guo2024dsagent,hong2025data,zhang2025data,sun2025survey}。许多系统可以在给定任务目标后完成若干分析步骤，却未必能够在面对具体数据集时形成与数据事实一致、与研究问题匹配、与后续执行资源协调的分析方案。对于结构化数据分析任务而言，这一差距尤为明显：如果规划阶段缺少对字段分布、跨列关系、图结构线索与样本语义特征的显式感知，系统就容易生成模板化、空泛化甚至与当前数据脱节的分析路径。

这一问题在专利分析场景中更加突出\citep{basberg1987patents,zhu2014patent,mejia2024patent}。专利数据同时具有技术属性、法律属性、时间属性和组织属性，研究问题往往要求系统同时处理变量解释、方法匹配、控制变量设置、图结构线索和时间演化等多个维度。现有方法若仅依赖问题文本和通用提示词，容易出现两类问题：一类是规划方案“看起来合理但并不贴合当前数据”，另一类是方案虽然具备局部亮点，但整体上缺乏方法组织和证据闭环，难以稳定形成可复验的分析链。

针对上述问题，本文关注一个更具体也更具可复用性的问题：在给定研究问题、结构化数据集和领域知识资源的前提下，如何构建一个能够先理解数据、再组织证据、最后输出可执行分析计划的多智能体规划框架。与仅优化执行器的路线不同\citep{yao2023react,wu2023autogen,hong2024metagpt,li2023camel}，本文把分析系统的关注点前移到规划阶段，并将其拆解为三个互补机制：其一，证据锚定负责在规划前提取字段统计、图结构特征和分层抽样语义样本，将“数据事实”显式注入上下文；其二，知识约束规划使用因果图谱和方法图谱分别约束假设空间与方法空间，减少自由生成带来的结构漂移；其三，迭代精化在首轮执行后识别证据缺口，并通过补证据、换方法、增控制和做稳健性分析的方式触发第二轮规划\citep{shinn2024reflexion,zhou2024hypothesis}。

本文进一步将期刊稿的论证目标从“内部模式比较”扩展为三层证据链：第一层通过四方案渐进与两组消融回答系统为何逐步变好；第二层通过 ReAct 风格单智能体基线和 execution-feedback 基线回答完整系统相较代表性外部机制路线的优势；第三层通过子领域稳定性验证、软件工程 Bugzilla 数据上的跨领域小规模验证，以及人工盲评评估方法的外部有效性与人类可接受性。换言之，本文不仅关注系统在单一实验设置中的得分高低，更关注规划框架是否能够在重复实验、外部对照和跨数据集场景下保持稳定优势。

本文的主要贡献如下。

\begin{itemize}
  \item 提出一种面向结构化分析任务的数据感知多智能体规划框架，将数据预览、图结构预览、知识约束和执行反馈统一纳入规划过程，并以结构化计划 JSON 作为统一中间表示层。
  \item 将知识约束进一步拆分为“假设空间约束”和“方法空间约束”两个层级，并通过双知识图谱实现协同约束，使系统能够显式区分“为何分析”和“如何分析”。
  \item 构建一套可复验的分层评估协议，将内部渐进实验、组件消融、文献启发的外部机制基线、子领域稳定性验证、跨领域小规模验证和人工盲评组织为统一证据链，用于评估框架的机制有效性与外部有效性。
  \item 在专利分析场景中验证该框架的工程可行性，并给出关于数据感知稳定收益、迭代机制上限作用以及知识图谱任务依赖性的经验结论；同时将软件工程数据上的跨领域验证纳入增强版投稿路线。
\end{itemize}

本文采用独立期刊稿结构，不复用学位论文“章节式系统说明”叙事。后续章节将依次给出问题定义、方法框架、系统实现、实验设计与结果讨论，并通过统一的统计与人工评审协议组织完整投稿证据链。
